\section{Introduction}\label{sec:intro}

Text artefacts that leave an organisation, such as an application log attached to a support ticket, a database export handed to an analyst, or a packet capture shared with a vendor, routinely carry two classes of sensitive value: personal data (e-mail addresses, telephone and card numbers, national identifiers, names) and credentials (application programming interface keys, access tokens, private keys, passwords). Data-loss-prevention systems inspect such content before it crosses a boundary \cite{alneyadi2016dlp}, and measurement studies show that credentials leak into shared artefacts at scale \cite{meli2019git,krause2023pushed}.

Two families of tools address the problem. Personal-data engines such as Microsoft Presidio \cite{presidio} combine pattern recognisers with a statistical named-entity model but ship no credential recognisers. Secret scanners such as gitleaks \cite{gitleaks} and TruffleHog \cite{trufflehog} target credentials in source repositories, and their accuracy varies widely across secret types \cite{basak2023comparative}. The artefacts that most need redaction are often too large to analyse in one call, so they are processed in pieces. A chunked redactor can then fail silently: a value that straddles a chunk boundary is presented to the detector in two halves, neither of which matches, and both halves are written out in the clear. The same problem is long established in network security, where signatures split across packets motivated stream reassembly in intrusion detection \cite{ptacek1998insertion} and streaming modes in regular-expression engines \cite{wang2019hyperscan,intel2019hyperscanintro}. In redaction, however, detectors read context around a match (``Patient ID:'' before a number, a key name before a token), confidences are adjusted by neighbouring text, and overlapping detections are merged before the text is rewritten. Overlapping windows are the obvious remedy, but how much overlap is enough depends on all of these, and the obvious answer, an overlap as long as the longest match, is wrong: a match whose right context is cut off at the window edge can be rejected there, so the boundary does not retreat out of it and its first part is emitted in the clear. None of the redaction or secret-scanning systems we examined states when its chunked output equals its whole-document output.

This paper answers that question for a concrete system, SafeStream-Redactor (hereafter SafeStream). Its contributions are the following.
\begin{enumerate}
\item \textbf{A tight, per-detector overlap condition for context-dependent redaction.} Every detector $d$ declares an upper bound $L_d$ on the span it can emit and bounds $C_{\mathrm{left},d}$ and $C_{\mathrm{right},d}$ on the context it reads (Assumptions~\ref{as:span} and~\ref{as:ctx}). We prove that an overlap $w \ge \max_d \max(L_d + C_{\mathrm{right},d},\, C_{\mathrm{left},d})$ makes the streamed output equal to the whole-document output for every chunking, and show that the term $L_d + C_{\mathrm{right},d}$ cannot be dropped: an overlap of $\max_d L_d$, the natural choice, can leak (Section~\ref{sec:formal}). We also measure the failure the condition prevents: on the test corpora, naive chunking leaked up to 237 gold values that a whole-document pass redacts, and overlapping windows still leaked some, while SafeStream's output was identical to the whole-document output under every chunking tried (Section~\ref{sec:equiv}). The underlying locality argument is standard; what is new is the condition for pipelines whose detectors read context on both sides, adjust confidences and resolve overlaps jointly, which pattern-matching engines that carry automaton state across blocks do not address. The condition is not tied to our detectors: wrapped around Presidio's recognisers, the engine made their chunked output identical to their whole-document output, and the same check exposed a Presidio pattern whose output depends on arbitrarily distant text (Section~\ref{sec:equiv}).
\item \textbf{A guarantee enforced by the code, not chosen by hand.} The engine computes the condition from the detectors' declarations and refuses to run with a smaller overlap; regression and differential tests check that the declarations hold. For the default pipeline the requirement is exactly the default overlap of 8192 characters.
\item \textbf{Parallel redaction that cannot lose or duplicate text.} A seam protocol between worker processes makes the parallel output equal to the sequential output under the same condition, and falls back to sequential processing on any disagreement, so no character is lost or emitted twice whatever the detectors do (Appendix~\ref{app:proofs}).
\item \textbf{Bounded memory, measured.} The working buffer is bounded independently of input size, and peak memory stays flat across a 256-fold range of input sizes, where whole-document redaction needs about 28 times the input size (Section~\ref{sec:memory}).
\item \textbf{A pre-registered detection evaluation on real data.} Equivalence says nothing about how good the detectors are. Beyond held-out test splits, we evaluate on four real, human-labelled corpora (system logs, 14,626 credentials from 337 repositories, court judgments and Wikipedia biographies), each with its protocol and manifest committed before any system was run, against Presidio and two secret scanners, with paired bootstrap intervals at the document, cluster and two-stage levels (Sections~\ref{sec:setup} and~\ref{sec:results}). SafeStream fully redacts more real credentials than gitleaks, robustly but with more false alarms; a free-text configuration, designed on development data and confirmed on 1,014 held-out judgments and on a second genre, detects more personal data than Presidio; and we report the margins that do not survive, including a reversal on court judgments for the configuration originally evaluated (Sections~\ref{sec:tab} and~\ref{sec:prose}).
\end{enumerate}

Three scope limits are stated at the outset. First, the guarantee concerns equivalence rather than accuracy: streaming does not improve detection, it makes the streamed result the same as the whole-document result, so every detection error of the whole-document pipeline is also an error of the streamed one. Second, the evaluation has known weaknesses that we quantify rather than omit: one corpus, RedactionBench, was inspected during rule development (Section~\ref{sec:contam}); two corpora are synthetic, and on real court judgments the evaluated named-entity configuration detected less personal data than Presidio until a free-text configuration was added (Section~\ref{sec:tab}); and fewer than one RedactionBench document in five was fully protected by any configuration we evaluated (Section~\ref{sec:noleak}). Third, SafeStream also ships an optional on-device neural tier and an interface for AI agents (a Model Context Protocol server and a REST API); their detection quality, what they disclose to an agent and their per-call cost are evaluated in a separate study \ssCiteAgent{} and are outside this paper. Every result here comes from the pattern tiers and, where stated, the spaCy tier; none depends on the neural tier or the agent interface. Section~\ref{sec:discussion} states when the approach is useful despite these limits and what evidence a deployment would need.
